\label{chap:apendice_exploratorios}

Antes do protocolo único, foram exploradas três arquiteturas KAN treinadas do zero, para ver até onde uma KAN chegaria sem um extrator de características pré-treinado. Elas ficam registradas aqui por completude. Os números vêm de outro cenário e \emph{não} são comparáveis com nenhum resultado do Capítulo~\ref{chap:results}.

\begin{itemize}
    \item \textbf{ResKANet-18:} uma rede residual cujas próprias convoluções são parametrizadas por splines (\texttt{reskanet\_18x32p} da \texttt{torch-conv-kan}~\cite{drokin2024torchconvkan}), treinada do zero com entradas de $128 \times 128$.
    \item \textbf{PureConvKAN:} três blocos convolucionais KAN (16, 32 e 64 canais), cada um seguido de max-pooling e normalização em lote, depois pooling médio global, Dropout(0,5) e uma camada linear de saída; treinada do zero com entradas de $128 \times 128$.
    \item \textbf{PureKAN:} a imagem achatada num vetor de pixels e passada diretamente a uma KAN, sem convolução; entradas de $64 \times 64$ ($3 \times 64 \times 64 = 12.288$ entradas).
\end{itemize}

A ResKANet-18 foi treinada só em entrada/saída, e a PureConvKAN e a PureKAN, nas duas tarefas. Todas foram treinadas por 60 épocas, de ponta a ponta, com AdamW e lotes de 8 imagens. Foram avaliadas num único hold-out 80/20 separado por caso, que evita o vazamento de casos, mas é uma única partição, com resolução menor e uma receita diferente da do protocolo único. A Tabela~\ref{tab:exploratorios} traz os resultados.

% Gerado por tcc_dissertation/scripts/gerar_tabelas.py a partir de resultados/. Não edite à mão.
\begin{table}[!htbp]
\centering
\caption{\EP{Exploratory KAN architectures trained from scratch (no ImageNet pre-training) in an earlier phase of this work. They used a single 80/20 hold-out split by case, inputs of 64 or 128 px, batches of 8 and their own recipe, so they are \emph{not} comparable with any number obtained under the single protocol. M = macro, m = micro, W = weighted.}{Arquiteturas KAN exploratórias treinadas do zero (sem pré-treino no ImageNet) numa fase anterior deste trabalho. Usaram um único hold-out 80/20 por caso, entradas de 64 ou 128 px, lotes de 8 e receita própria, portanto \emph{não} são comparáveis com nenhum número obtido no protocolo único. M = macro, m = micro, W = ponderada.}}
\label{tab:exploratorios}
\footnotesize
\setlength{\tabcolsep}{2.5pt}
\resizebox{\textwidth}{!}{%
\begin{tabular}{llcccccccccccccccc}
\toprule
 & & Acc & \multicolumn{3}{c}{\EP{Precision}{Precisão}} & \multicolumn{3}{c}{\EP{Recall}{Revocação}} & \multicolumn{3}{c}{F1} & \multicolumn{3}{c}{AUC} & \multicolumn{3}{c}{\EP{Specificity}{Especificidade}} \\
\EP{Model}{Modelo} & \EP{Input}{Entrada} & & M & m & W & M & m & W & M & m & W & M & m & W & M & m & W \\
\midrule
\multicolumn{18}{l}{\textit{\EP{Entry/exit}{Entrada/saída}}} \\
ResKANet-18 & 128 px & 0.822 & 0.767 & 0.822 & 0.820 & 0.762 & 0.822 & 0.822 & 0.765 & 0.822 & 0.821 & 0.839 & 0.885 & 0.839 & 0.762 & 0.822 & 0.703 \\
PureConvKAN & 128 px & 0.707 & 0.614 & 0.707 & 0.704 & 0.611 & 0.707 & 0.707 & 0.613 & 0.707 & 0.706 & 0.682 & 0.768 & 0.682 & 0.611 & 0.707 & 0.516 \\
PureKAN & 64 px & 0.692 & 0.614 & 0.692 & 0.692 & 0.614 & 0.692 & 0.692 & 0.614 & 0.692 & 0.692 & 0.684 & 0.733 & 0.684 & 0.614 & 0.692 & 0.536 \\
\midrule
\multicolumn{18}{l}{\textit{\EP{MLSD}{MLSD}}} \\
PureConvKAN & 128 px & 0.803 & 0.268 & 0.803 & 0.645 & 0.333 & 0.803 & 0.803 & 0.297 & 0.803 & 0.716 & 0.536 & 0.888 & 0.509 & 0.667 & 0.902 & 0.197 \\
PureKAN & 64 px & 0.712 & 0.355 & 0.712 & 0.689 & 0.324 & 0.712 & 0.712 & 0.328 & 0.712 & 0.693 & 0.463 & 0.841 & 0.496 & 0.679 & 0.856 & 0.326 \\
\bottomrule
\end{tabular}}
\end{table}


A única leitura que esses números sustentam é qualitativa: sem um extrator pré-treinado, as arquiteturas KAN chegam a 0,69--0,82 de acurácia em entrada/saída, e na tarefa de distância a AUC macro das duas fica perto do acaso ou abaixo dele (0,46--0,54), ou seja, aprendem pouco sobre as classes minoritárias. Isso é coerente com o achado principal de que é o backbone pré-treinado, e não a KAN, que faz a discriminação.
